\chapter{Introdução}
\textbf{Utilize este modelo para construir seu projeto de pesquisa. O livro de \textcite{jalongo2016writing} foi utilizado para compor esse material e lá você pode encontrar mais informações sobre como escrever para publicação.}

\textbf{Comece apresentando uma introdução sobre o tema proposto: principais conceitos e aplicações.}

\lipsum[1-1]

\section{Contextualização}
\textbf{Qual a importância do tema? Quais os fatores que determinam a escolha do tema? Qual a sua relação com o tema? Qual a sua contribuição para o avanço desse tema?}

\lipsum[1-2]

\section{Objetivo}

\textbf{Deixe claro qual o problema que você pretende responder com a pesquisa, assim como sua delimitação espacial e temporal. Apresente os objetivos em termos claros e precisos, listando possíveis sub-problemas. Utilize verbos de ação como verificar, identificar, descrever e analisar.}

\textbf{Se houver hipótese, deixar explicitas as relações da hipotes com as variáveis do problema}

\lipsum[1-2]
